% Generated draft: requires mathematical and source review.
\documentclass[11pt]{article}
\usepackage[margin=25mm]{geometry}
\usepackage{fontspec}
\setmainfont{Latin Modern Roman}
\usepackage{amsmath,amssymb,mathtools,unicode-math}
\setmathfont{Latin Modern Math}
\usepackage{enumitem,needspace,xurl,hyperref}
\setlength{\parindent}{0pt}
\setlength{\parskip}{5pt}
\setlength{\emergencystretch}{3em}
\newcommand{\N}{\mathbb{N}}
\newcommand{\Z}{\mathbb{Z}}
\newcommand{\Q}{\mathbb{Q}}
\newcommand{\R}{\mathbb{R}}
\newcommand{\C}{\mathbb{C}}
\newcommand{\F}{\mathbb{F}}
\newcommand{\A}{\mathbb{A}}
\newcommand{\PP}{\mathbb{P}}
\newcommand{\E}{\mathbb{E}}
\newcommand{\eps}{\varepsilon}
\newcommand{\dd}{\,\mathrm{d}}
\begin{document}
\section*{UnsolvedMath ID: 30006339}
\textbf{Source problem code:} catalogue entry
\subsection{Definitions and mathematical statement}
Let \(n\geq 1\) denote the rank. In the setting of the cited construction, let \(X\) be the underlying space and let \(C\) be its \(p\)-adic coefficient field. The supplied record does not specify the hypotheses on \(X\), the precise field \(C\), or the topology and category of the moduli spaces; these remain qualifications on the formulation below.

Write \(\Omega_X\) for the differential sheaf occurring in the source, and write \((-i)\) for its indicated Tate twist. The source describes the Hitchin base as the affine vector space
\[
 A_{n,X}=\bigoplus_{i=1}^{n}\left(H^0\!\left(X,\Omega_X^{\otimes i}(-i)\right)\otimes_C\mathbb A_C^1\right),
\]
where tensoring with \(\mathbb A_C^1\) denotes the affine vector space associated with the indicated \(C\)-vector space. Its coefficient tuples are viewed as monic polynomials
\[
 f(T)=T^n+a_{n-1}T^{n-1}+\cdots+a_0,
 \qquad
 a_{n-i}\in H^0\!\left(X,\Omega_X^{\otimes i}(-i)\right)
 \quad(1\leq i\leq n).
\]
Here these polynomial coefficients are sections of different tensor powers, with multiplication interpreted accordingly.

The source describes \(\operatorname{Rep}_{n,X}\) as the \(p\)-adic character variety parametrising irreducible continuous \(n\)-dimensional \(C\)-representations of \(\pi_1(X)\). The precise fundamental group, topology, and moduli equivalence are not supplied. It introduces a \(p\)-adic Hitchin morphism
\[
 \widetilde H_n:\operatorname{Rep}_{n,X}\longrightarrow A_{n,X},
\]
obtained through moduli of \(v\)-vector bundles and described as a twist of the characteristic-polynomial map for Higgs fields.

\textbf{Question.} For the ranks \(n\) and spaces \(X\) covered by the cited construction, is \(\widetilde H_n\) surjective? The intended range of \(X\) and the meaning of surjectivity in the relevant category must be recovered from the cited construction. In particular, the supplied record does not justify replacing this question by surjectivity on \(C\)-valued points.

\textbf{Supplied qualifications and context.} The excerpt reports that \(\widetilde H_n\) is open and equivariant for a natural Tate-twisted semilinear \(G_K\)-action, while surjectivity is conjectural. Neither \(K\) nor this action is defined in the supplied record. It also reports compatibility with direct sums and rank-one twisting:
\[
 \widetilde H_{n+m}(V\oplus W)
 =\widetilde H_n(V)\,\widetilde H_m(W),
\]
\[
 \widetilde H_n(L\otimes V)(T)
 =\widetilde H_n(V)\bigl(T-\widetilde H_1(L)\bigr),
\]
where \(L\) has rank one and the first product is multiplication of the corresponding polynomials. These are reported contextual properties, not additional conclusions of this draft. The direct-sum formula requires clarification: a direct sum of two positive-rank representations is reducible and therefore does not belong to a space parametrising only irreducible representations. The excerpt must be using a broader representation space for that formula, or omitting a qualification; this draft does not choose between those possibilities.

The background attributes the construction to Faltings' \(p\)-adic Simpson correspondence, presented there as a comparison between \(v\)-vector bundles and Higgs bundles. It reports joint work with Daxin Xu identifying the moduli of stable \(v\)-vector bundles with a twist of the moduli of stable Higgs bundles over the Hitchin base, and a natural open embedding of the irreducible character variety into the moduli of rank-\(n\) \(v\)-vector bundles. The hypotheses and precise meanings of these comparisons are absent from the supplied excerpt and are not asserted here without qualification.
\subsection{Short English statement}
The question asks whether every point of the intended Hitchin base is attained by the \(p\)-adic Hitchin morphism from the specified character variety. The supplied source calls this conjectural and reports openness and equivariance. A complete formulation still requires the hypotheses on \(X\), the coefficient field, the moduli category, and the intended meaning of surjectivity. The excerpt also leaves an incompatibility between its irreducible-representation description and its direct-sum formula.
\subsection{Sources}
\begin{itemize}
\item Local UnsolvedMath registry, supplied complete source record: permanent ID \(30006339\), exact source problem code \texttt{catalogue entry}, title ``Surjectivity of the p-Adic Hitchin Morphism.'' The record assigns the problem to the Oberwolfach Reports Open Problems collection and gives proposed year 2025; no individual proposer is supplied.
\item \emph{Non-Archimedean Geometry and Applications} (2025), Oberwolfach Reports, pp.\ 1517--1578, DOI \texttt{10.4171/owr/2025/29}. Supplied report reference and literature-source link: \url{https://doi.org/10.4171/owr/2025/29}. No precise page or contribution author for this extracted question is supplied.
\item \emph{p-adic non-abelian Hodge theory for curves}. Supplied literature-source link: \url{https://arxiv.org/abs/2402.01365}.
\item \emph{Moduli spaces in p-adic non-abelian Hodge theory}. Supplied literature-source link: \url{https://arxiv.org/abs/2207.13819}.
\item Supplied literature assessment, checked on \texttt{2026-08-22}: the morphism is reported proved open and equivariant, with general surjectivity explicitly conjectural; the supplied curation reports finding no later primary source resolving general surjectivity by that cutoff. This is supplied dated curation derived from the OWR record, not an independently verified literature assessment. The supplied \texttt{open} status and \texttt{corrected\_verified} statement label are registry labels; the latter is described as verification of formula delimiters, notation, grammar, and typography while preserving mathematical scope.
\end{itemize}
\end{document}
